\documentclass[10pt,a4paper]{amsart}
\usepackage[margin=19mm]{geometry}
\usepackage{amsmath,amssymb,mathtools}
\usepackage[scr=rsfs,cal=euler]{mathalfa}
\usepackage{xfrac}
\usepackage[unicode,hidelinks,bookmarks=false]{hyperref}
\newcommand{\ie}{i.e.\ }
\newcommand{\cf}{cf.\ }
\newcommand{\resp}{respectively\ }
\newcommand{\Subsection}{\subsection}
\providecommand{\nobreakdash}{\nobreak\hbox{-}}
\setlength{\emergencystretch}{2em}
\multlinegap0pt
\usepackage{amssymb}
\newtheorem{lemma}{Lemma}
\newtheorem{theorem}{Theorem}
\title{Semigroups, Cartier divisors and convex bodies}
\author{Askold Khovanskii}
\date{Source: arXiv:2502.13099v2 (CC BY 4.0)}
\begin{document}
\maketitle

\noindent\emph{Source and licence.} Semigroups, Cartier divisors and convex bodies. Version arXiv:2502.13099v2 (CC BY 4.0), distributed by arXiv under the Creative Commons Attribution 4.0 International licence, http://creativecommons.org/licenses/by/4.0/, as recorded in the arXiv licence metadata for this exact version and shown on the abstract page https://arxiv.org/abs/2502.13099v2. Full licence notice: \texttt{licenses/arxiv-2502.13099-CC-BY-4.0.txt}.\par\smallskip

\noindent\textit{Editorial scope.} Counted unit: the article's theorem strongKht (source0.tex lines 1051--1056). The article's complete original proof of it is delivered with it (source0.tex lines 1058--1073, one \texttt{proof} environment of 16 lines). No mathematics is written, repaired or re-derived: the statement and the proof are the article's own lines, in the article's own notation, and no retained line is substituted or re-flowed.  Retained uncounted as the source-local context the proof needs: lines 1006--1008.  Nothing was omitted from any retained span, so the package carries no omission record.  No retained line contains a citation, so the wrapper carries no bibliography and no bibliography entry.  No retained line contains a figure environment, an image-inclusion command or drawing code, so no image is delivered and the package declares no asset dependency.  Editorial formatting only, in addition: an AMS \texttt{amsart} preamble that restates the article's own declarations and macros used by the retained lines, namely the environments \texttt{lemma}, \texttt{theorem} on separate counters, together with the standard packages those lines need.  The retained environments print in this document as Lemma 1, Theorem 1 in order of appearance, whereas the article prints the same items with the numbering of its own full document.  The complete native source of the article as distributed in the arXiv e-print for this exact version is bundled unchanged as \texttt{sources/173500/source0.tex}.
\begin{lemma}[Algebraic version of Alexandrov--Fenchel  inequality]\label{weakKht} If $L_3,\dots,L_n$ are very big  subspaces, then
\begin{equation*}[L_1, L_2, L_3,\dots, L_n]^2 \geq [L_1,L_1, L_3, \dots, L_n][L_2,L_2, L_3, \dots, L_n].\end{equation*}
\end{lemma}

\begin{theorem}\label{strongKht} If $\mathcal L_1, \mathcal L_2, \mathcal L_3,\dots,\mathcal L_n$ are
psef type elements in the group $G(X)$, then
\begin{equation}\label{nefaf}
[\mathcal L_1, \mathcal L_2,\mathcal  L_3,\dots, \mathcal L_n]^2 \geq [\mathcal L_1, \mathcal L_1, \mathcal  L_3, \dots,\mathcal L_n] [ \mathcal L_2, \mathcal L_2,   \mathcal L_3, \dots, \mathcal  L_n].
\end{equation}
\end{theorem}

\begin{proof} If each  $\mathcal L_i$ is a very big element in $G(X)$, then theorem follows from Lemma \ref{weakKht}, since the intersection index respects the equivalence $\sim$ in $K(X)$.

Assume that $\mathcal L_1,\dots,\mathcal L_n$ are big elements, i.e., for some natural number $k_1,\dots,k_n$ the elements $\mathcal L_1^{k_1},\dots,\mathcal L_n^{k_n}$ are very big  elements. Put $k=k_1\cdot \ldots \cdot k_n$. Then elements $\mathcal L_1^k,\dots,\mathcal L_n^k$ are very big spaces. Thus, for the $n$-tuple of spaces $\mathcal L_1^k,\dots,\mathcal L_n^k$  inequality (\ref{nefaf}) holds. It implies inequality (\ref{nefaf}) for the $n$-tuple $\mathcal L_1,\dots,\mathcal L_n$, since the intersection index is multi-linear.


Assume that  $\mathcal L_i$ are  psef type elements, i.e, there are  big elements $\mathcal L_i'$ such that for any natural number $q$ all elements $\mathcal L_i^q\mathcal L_i'$ are  big. Put $\mathcal L=\mathcal L_1'\cdot \ldots \cdot \mathcal L_n'$. Then all elements $\mathcal L_i^q\mathcal L$ are  big.
Instead of elements
$\mathcal L_1,...,\mathcal L_n$,
let us consider the elements  $\mathcal L_1(q) =\mathcal L_1^q \mathcal L$ $,\dots$,   $\mathcal L_n(q) =\mathcal L_n^q \mathcal L$
for every  natural number  $q$. For the element $\mathcal L_1(q),\dots, \mathcal L_n(q)$ the Alexandrov--Fenchel type inequality holds, i.e., for every natural number $q$, we  have $A(q)\geq B(q)$, where $A(q)$ and $B(q)$ are the following polynomials in $q$:
\begin{equation*}A(q)=[\mathcal L_1(q),\mathcal L_2(q),\mathcal  L_3(q),\dots,\mathcal  L_n(q)]^2; \end{equation*}  \begin{equation*}B(q)= [\mathcal L_1(q), \mathcal L_1(q), \mathcal  L_3(q), \dots,\mathcal L_n(q)][\mathcal L_2(q), \mathcal L_2(q), \mathcal L_3 (q), \dots, \mathcal  L_n (q)].\end{equation*}


The polynomials $A(q)$ and $B(q)$ have the following  leading terms  $A_{2n}$, $B_{2n}$:  \begin{equation*}A_{2n}=q^{2n}  [\mathcal L_1, \mathcal L_2, \mathcal L_3,\dots, \mathcal L_n]^2;\end{equation*}
\begin{equation*}B_{2n}=q^{2n} [\mathcal L_1,\mathcal L_1, \mathcal  L_3, \dots,\mathcal L_n] [\mathcal L_2,\mathcal L_2, \mathcal  L_3, \dots, \mathcal  L_n].\end{equation*}  The inequalities $A(q)\geq B(q)$  imply that $A_{2n}\geq B_{2n}$ which proves inequality (\ref{nefaf}).
\end{proof}

\end{document}
